% ══════════════════════════════════════════════════════════════════
%  BACKGROUNDS — full-bleed decorative TikZ backgrounds.
%  Each macro takes ONE argument: opacity (0–1).
%  Use inside a {tikzpicture}[remember picture, overlay] anchored to
%  current page, or call \FullBleed{<color>}{<macro>{opacity}} below.
% ══════════════════════════════════════════════════════════════════

% Paint the whole page a color, then draw a pattern over it.
% Usage: \FullBleed{MidnightNavy}{\constellation{0.35}}
\newcommand{\FullBleed}[2]{%
  \begin{tikzpicture}[remember picture, overlay]
    \fill[#1] (current page.south west) rectangle (current page.north east);
    \begin{scope}[shift={(current page.south west)}]
      #2
    \end{scope}
  \end{tikzpicture}%
}

% ── \circuitboard{opacity} — PCB traces with pads ──
\newcommand{\circuitboard}[1]{%
  \begin{scope}[opacity=#1, NeonCyan]
    \foreach \i in {1,...,18}{
      \pgfmathsetmacro{\sx}{mod(\i*47, 200)/10}
      \pgfmathsetmacro{\sy}{mod(\i*89, 280)/10}
      \pgfmathsetmacro{\la}{1+mod(\i*13,4)}
      \pgfmathsetmacro{\lb}{1+mod(\i*7,5)}
      \draw[line width=0.5pt] (\sx,\sy) -- ++(\la,0) -- ++(0,\lb)
        -- ++({mod(\i,2)==0 ? 1.4 : -1.4},0);
      \fill (\sx,\sy) circle (1.6pt);
      \draw ({\sx+\la+(mod(\i,2)==0 ? 1.4 : -1.4)},{\sy+\lb}) circle (1.9pt);
    }
  \end{scope}%
}

% ── \constellation{opacity} — dot grid + connecting lines ──
\newcommand{\constellation}[1]{%
  \begin{scope}[opacity=#1]
    \foreach \i in {1,...,60}{
      \pgfmathsetmacro{\sx}{mod(\i*73, 210)/10}
      \pgfmathsetmacro{\sy}{mod(\i*131, 297)/10}
      \pgfmathsetmacro{\r}{0.5+0.7*rnd}
      \fill[IceWhite] (\sx,\sy) circle (\r pt);
    }
    \foreach \a/\b/\c/\d in {2.1/24.3/5.8/22.1, 5.8/22.1/9.2/25.4,
      12.4/8.3/15.1/11.2, 15.1/11.2/18.3/9.7, 3.4/13.1/6.7/15.8,
      16.9/19.4/19.5/17.2, 8.8/3.9/11.6/5.4}{
      \draw[IceWhite, line width=0.35pt] (\a,\b) -- (\c,\d);
    }
  \end{scope}%
}

% ── \topography{opacity} — organic contour lines ──
\newcommand{\topography}[1]{%
  \begin{scope}[opacity=#1]
    \foreach \i in {0,...,11}{
      \pgfmathsetmacro{\yy}{2.2*\i}
      \draw[FogGray, line width=0.45pt]
        (-1,\yy) .. controls (4,{\yy+1.6+0.9*sin(\i*70)}) and (8,{\yy-1.4})
        .. (12,{\yy+1.1*cos(\i*45)}) .. controls (16,{\yy+2.1}) and (19,{\yy-1.2})
        .. (22,{\yy+0.6});
    }
  \end{scope}%
}

% ── \hexgrid{opacity} — hexagonal grid ──
\newcommand{\hexgrid}[1]{%
  \begin{scope}[opacity=#1]
    \foreach \row in {0,...,13}{
      \foreach \col in {0,...,8}{
        \pgfmathsetmacro{\hx}{\col*2.6 + mod(\row,2)*1.3}
        \pgfmathsetmacro{\hy}{\row*2.25}
        \draw[StormBlue!70!IceWhite, line width=0.4pt]
          ({\hx+1.5*cos(30)},{\hy+1.5*sin(30)})
          \foreach \ang in {90,150,210,270,330}{
            -- ({\hx+1.5*cos(\ang)},{\hy+1.5*sin(\ang)})
          } -- cycle;
      }
    }
  \end{scope}%
}

% ── \stardust{opacity} — scattered dots and rings ──
\newcommand{\stardust}[1]{%
  \begin{scope}[opacity=#1]
    \foreach \i in {1,...,80}{
      \pgfmathsetmacro{\sx}{21*rnd}
      \pgfmathsetmacro{\sy}{29.7*rnd}
      \pgfmathsetmacro{\r}{0.3+0.8*rnd}
      \fill[IceWhite] (\sx,\sy) circle (\r pt);
    }
    \foreach \i in {1,...,7}{
      \pgfmathsetmacro{\sx}{21*rnd}
      \pgfmathsetmacro{\sy}{29.7*rnd}
      \draw[IceWhite, line width=0.3pt] (\sx,\sy) circle ({(2+4*rnd)*1pt});
    }
  \end{scope}%
}

% ── \blueprint{opacity} — architectural grid + axis marks ──
\newcommand{\blueprint}[1]{%
  \begin{scope}[opacity=#1]
    \draw[IceWhite, line width=0.18pt, step=0.5] (0,0) grid (21,29.7);
    \draw[IceWhite, line width=0.5pt, step=2.5] (0,0) grid (21,29.7);
    \foreach \x [count=\i] in {2.5,5,...,20}{
      \node[IceWhite, font=\ttfamily\fontsize{5pt}{5pt}\selectfont]
        at (\x,29.3) {\i};
    }
    \foreach \y/\lab in {2.5/A,5/B,7.5/C,10/D,12.5/E,15/F,17.5/G,20/H,22.5/I,25/J,27.5/K}{
      \node[IceWhite, font=\ttfamily\fontsize{5pt}{5pt}\selectfont]
        at (0.35,\y) {\lab};
    }
  \end{scope}%
}

% ── \ledgerlines{opacity} — Verita's own ruled paper (for artifacts) ──
\newcommand{\ledgerlines}[1]{%
  \begin{scope}[opacity=#1]
    \foreach \y in {0.8,1.6,...,29.2}{
      \draw[IronGallInk!25, line width=0.3pt] (0,\y) -- (21,\y);
    }
  \end{scope}%
}
